\documentclass[14pt]{extarticle}
\usepackage{fontspec}
\usepackage{polyglossia}
\usepackage{amsmath}
\usepackage{amssymb}
\usepackage{xcolor}
\usepackage{fancyhdr}
\usepackage{graphicx}
\usepackage{listings}
\usepackage[most]{tcolorbox}
% \usepackage{setspace}
\usepackage{pifont}
\usepackage{enumitem}
\usepackage{titlesec}
\usepackage[bottom]{footmisc}
\usepackage{titling}
\usepackage{minted}
\usepackage{etoolbox}
\usepackage{array}
\usepackage{longtable} % page-breakable tables (must load before polyglossia/bidi RTL setup)
\usepackage{extsizes}
\usepackage{float}
\usepackage{booktabs}
\usepackage{hyperref}
\usepackage[onehalfspacing]{setspace}
\newif\ifwhiteblack
\ifdefined\setwhiteblack
  \setwhiteblack
\fi

\ifwhiteblack
  \usepackage[a4paper, landscape, top=1.25cm, bottom=1.25cm, left=1.25cm, right=1.25cm, headheight=1.5cm, headsep=0.5cm]{geometry}
\else
  \usepackage[a4paper, top=2.7cm, bottom=2.7cm, left=2.5cm, right=2.5cm, headheight=1.5cm, headsep=0.5cm]{geometry}
\fi
\usepackage{tikz}
\usepackage{datetime2}

\usetikzlibrary{automata, positioning, arrows.meta, arrows}
\providecommand{\LTcaptype}{table}

\newfontfamily\emoji{Segoe UI Emoji}

\pagestyle{fancy}

\setmainlanguage[numerals=western]{arabic}
\setotherlanguage{english}
\newfontfamily\arabicfont[Script=Arabic]{Calibri}
\newfontfamily\arabicfonttt[Script=Arabic]{Courier New}



\ifwhiteblack
\lstset{
  language=[Sharp]C,
  numbers=left,
  stepnumber=1,
  numbersep=8pt,
  frame=single,
  basicstyle=\ttfamily\small,
  keywordstyle=\bfseries,
  stringstyle=\ttfamily,
  commentstyle=\itshape
}
\else
\lstset{
  language=[Sharp]C,
  numbers=left,
  stepnumber=1,
  numbersep=8pt,
  frame=single,
  basicstyle=\ttfamily\small,
  keywordstyle=\color{blue},
  stringstyle=\color{red},
  commentstyle=\color{green!50!black}
}
\fi

\newif\ifdetailed
\ifdefined\setdetailed
  \setdetailed
\fi

\newif\ifwithsols
\ifdefined\setwithsols
  \setwithsols
\fi

\ifwhiteblack
  \AtBeginDocument{
    \LARGE
    \usemintedstyle{bw}
  }
\fi

% Default minted look: compact code that wraps long lines.
% Not applied in black-and-white landscape mode, which already uses \LARGE.
\ifwhiteblack\else
  \setminted{fontsize=\small, breaklines=true}
\fi

% unified tcolorboxes for articles
\ifwhiteblack
\tcbset{colback=white, colframe=black, fonttitle=\bfseries, boxrule=1pt}
\newtcolorbox{boxDef}[1][]{title={{\emoji📘} تعريف\ifx\\#1\\\else ~#1\fi :}}
\newtcolorbox{boxExercise}[1][]{title={{\emoji🧩} تمرين\ifx\\#1\\\else ~#1\fi :}}
\newtcolorbox{boxExample}[1][]{title={{\emoji📝} مثال\ifx\\#1\\\else ~#1\fi :}}
\newtcolorbox{boxNote}[1][]{title={{\emoji✨} ملاحظة\ifx\\#1\\\else ~#1\fi :}}
\newtcolorbox{boxAttention}[1][]{title={{\emoji🔔} تنبيه\ifx\\#1\\\else ~#1\fi :}}
\newtcolorbox{boxWarning}[1][]{title={{\emoji⚡} ملاحظة هامة\ifx\\#1\\\else ~#1\fi :}}
\newtcolorbox{boxSolution}[1][]{title={{\emoji✅} حل\ifx\\#1\\\else ~#1\fi :}}
\newtcolorbox{boxSymbol}[1][]{title={{\emoji🔣} رمز\ifx\\#1\\\else ~#1\fi :}}
\newtcolorbox{boxHint}[1][]{title={{\emoji💡} تلميح\ifx\\#1\\\else ~#1\fi :}}
\else
\tcbset{colback=white, colframe=black, fonttitle=\bfseries, boxrule=0.8pt}
\newtcolorbox{boxDef}[1][]{colback=blue!5!white,colframe=blue!75!black,
  title={{\emoji📘} تعريف\ifx\\#1\\\else ~#1\fi :}}
\newtcolorbox{boxExercise}[1][]{colback=cyan!5!white,colframe=cyan!70!black,
  title={{\emoji🧩} تمرين\ifx\\#1\\\else ~#1\fi :}}
\newtcolorbox{boxExample}[1][]{colback=yellow!5!white,colframe=orange!90!black,
  title={{\emoji📝} مثال\ifx\\#1\\\else ~#1\fi :}}
\newtcolorbox{boxNote}[1][]{colback=gray!10!white,colframe=black,
  title={{\emoji✨} ملاحظة\ifx\\#1\\\else ~#1\fi :}}
\newtcolorbox{boxAttention}[1][]{colback=magenta!10!white,colframe=magenta!80!black,
  title={{\emoji🔔} تنبيه\ifx\\#1\\\else ~#1\fi :}}
\newtcolorbox{boxWarning}[1][]{colback=red!5!white,colframe=red!75!black,
  title={{\emoji⚡} ملاحظة هامة\ifx\\#1\\\else ~#1\fi :}}
\newtcolorbox{boxSolution}[1][]{colback=green!5!white,colframe=green!60!black,
  title={{\emoji✅} حل\ifx\\#1\\\else ~#1\fi :}}
\newtcolorbox{boxSymbol}[1][]{colback=purple!5!white,colframe=purple!70!black,
  title={{\emoji🔣} رمز\ifx\\#1\\\else ~#1\fi :}}
\newtcolorbox{boxHint}[1][]{colback=teal!5!white,colframe=teal!60!black,
  title={{\emoji💡} تلميح\ifx\\#1\\\else ~#1\fi :}}
\fi


\ifwhiteblack
\tcbset{simplecode/.style={ colback=white, colframe=black, boxrule=1pt, arc=2pt, left=4pt,right=4pt,top=4pt,bottom=4pt}}
\else
\tcbset{simplecode/.style={ colback=gray!5, colframe=black!50, boxrule=0.4pt, arc=2pt, left=4pt,right=4pt,top=4pt,bottom=4pt}}
\fi
% Compact code boxes: tighter line spacing (the preamble uses onehalfspacing) and small vertical gaps.
% Must come AFTER the simplecode definition above.
\tcbset{simplecode/.append style={fontupper=\setstretch{0.95}, top=4pt, bottom=4pt, before skip=6pt, after skip=6pt}}
\newenvironment{boxCode}{\begin{tcolorbox}[simplecode]}{\end{tcolorbox}}

% Code cell used inside tables (keeps code left-to-right inside RTL tables).
% Use inside a fixed-width p{...} column: \begin{pycell}\begin{minted}...\end{minted}\end{pycell}
\newtcolorbox{pycell}{frame empty,colback=white,boxsep=0pt,left=2pt,right=2pt,top=2pt,bottom=2pt,width=\linewidth}

\newcolumntype{C}[1]{>{\centering\arraybackslash}p{#1}}

\makeatletter
\DTMnewdatestyle{dotted}{%
  \renewcommand{\DTMdisplaydate}[4]{\two@digits{##3}.\two@digits{##2}.##1}%
  \renewcommand{\DTMDisplaydate}{\DTMdisplaydate}%
}
\makeatother
\DTMsetdatestyle{dotted}

% redefine spaces after titles
\makeatletter
\renewcommand{\@maketitle}{%
  \begin{center}
    {\huge \bfseries \@title \par}%
    \vskip 0.2em % space between title and author
    {\large \@author \par}%
    % \vskip 0.2em % space between author and date
    % {\normalsize \@date \par}%
    {\small \hfill \begin{english}\DTMtoday\end{english} \par}%
  \end{center}
}
\makeatother

\fancyhf{} % clear default
\fancypagestyle{plain}{
  \fancyhf{}
  \fancyhead[L]{مدرسة التسامح الشاملة}
  % \fancyhead[L]{\includegraphics[height=1cm]{../../../images/logoTasamoh.png}}
  \fancyhead[R]{الأستاذ محمود اغبارية}
  \fancyfoot[C]{\thepage}
}

\fancyhead[L]{مدرسة التسامح الشاملة}
\fancyhead[R]{الأستاذ محمود اغبارية}
\fancyfoot[C]{\thepage}

\setcounter{tocdepth}{3} % only section subsection and subsubsection in TOC

\newcommand{\importpdfpage}[2]{%
    \thispagestyle{empty}
    \newgeometry{left=0cm,right=0cm,top=0cm,bottom=0cm}
    \includegraphics[page=#2,width=0.95\paperwidth,height=0.95\paperheight,keepaspectratio]{#1}
    \restoregeometry
}

\newcommand{\insertFullImg}[1]{%
    \noindent
    \makebox[\textwidth][c]{
        \includegraphics[width=0.9\paperwidth,keepaspectratio]{#1}%
    }%
}%


\ifwithsols
\title{حل ورقة عمل 1 - الفئة \textenglish{Point}}
\else
\title{ورقة عمل 1 - الفئة \textenglish{Point}}
\fi

\begin{document}

\maketitle


أمامك واجهة الفئة \textenglish{Point} التي تمثّل نقطة في هيئة المحاور:

\begin{center}
\renewcommand{\arraystretch}{1.8}
{\small
\begin{tabular}{|p{0.5\textwidth}|l|}
\hline
\multicolumn{2}{|c|}{\textbf{\textenglish{Point}}} \\
\hline
صفة تُمثّل الإحداثي $x$ للنقطة في المستوى. & \textenglish{private double x} \\
\hline
صفة تُمثّل الإحداثي $y$ للنقطة في المستوى. & \textenglish{private double y} \\
\hline
عمليّة بنائيّة تتلقّى قيمًا لجميع صفات الكائن. & \textenglish{public Point(double x, double y)} \\
\hline
عمليّة بنائيّة (افتراضية) تُهيّئ النقطة إلى نقطة الأصل $(0,0)$. & \textenglish{public Point()} \\
\hline
عمليّة بنائيّة تتلقّى كائنًا آخر من نمط \textenglish{Point} وتنسخ قيمه (Copy Constructor). & \textenglish{public Point(Point other)} \\
\hline
عمليّة تُعيد قيمة الإحداثي $x$. & \textenglish{public double GetX()} \\
\hline
عمليّة تُعيد قيمة الإحداثي $y$. & \textenglish{public double GetY()} \\
\hline
عمليّة تُعدّل قيمة الإحداثي $x$ للنقطة الحالية. & \textenglish{public void SetX(double x)} \\
\hline
عمليّة تُعدّل قيمة الإحداثي $y$ للنقطة الحالية. & \textenglish{public void SetY(double y)} \\
\hline
عمليّة تتلقّى قيمتين وتُزيح (تُحرّك) النقطة الحالية بمقدار $x$ على محور السينات وبمقدار $y$ على محور الصادات. & \textenglish{public void MoveBy(double x, double y)} \\
\hline
عمليّة تُعيد نصّاً يُمثّل النقطة بالصيغة $(x, y)$. & \textenglish{public string ToString()} \\
\hline
\end{tabular}
}
\end{center}

\clearpage

\begin{enumerate}[label=\textbf{\arabic*.}]
    \item \textbf{تطبيق الفئة:} \\
    اكتب تطبيقًا كاملاً للفئة \textenglish{Point} يشمل تعريف الخصائص، العمليات البنائية الثلاث، عمليات الاسترجاع والتعديل، وعمليات \textenglish{MoveBy} و \textenglish{ToString}.

    \item \textbf{استخدام الفئة في البرنامج الرئيسي:} \\
    اكتب برنامجًا رئيسيًا (\textenglish{Main}) يقوم بإنشاء نقطتين مختلفتين، ويطبع إحداثياتهما، ثم يفحص الدوال التي قمت بكتابتها في الفئة.

    \item \textbf{دوال ثابتة (\textenglish{Static Methods}) داخل الفئة \textenglish{Point}:} \\
    أضف الدوال الثابتة التالية إلى الفئة \textenglish{Point}، ثم استدعها في البرنامج الرئيسي لاختبارها:
    \begin{enumerate}[label=\textbf{\alph*)}]
        \item عمليّة باسم \textenglish{Slope} تتلقّى كائنين من نمط \textenglish{Point} وتُعيد ميل المستقيم المار بينهما. (قانون الميل: $m = \frac{y_2 - y_1}{x_2 - x_1}$)
        \item عمليّة باسم \textenglish{Distance} تتلقّى كائنين من نمط \textenglish{Point} وتُعيد المسافة بينهما. (قانون المسافة: $d = \sqrt{(x_1 - x_2)^2 + (y_1 - y_2)^2}$)
        \item عمليّة باسم \textenglish{GetMiddlePoint} تتلقّى كائنين من نمط \textenglish{Point} وتُعيد كائنًا جديدًا يُمثّل نقطة المنتصف بينهما. (قانون نقطة المنتصف: $M = (\frac{x_1 + x_2}{2}, \frac{y_1 + y_2}{2})$)
    \end{enumerate}

    \item \textbf{دوال مساعدة في البرنامج الرئيسي (خارج الفئة \textenglish{Point}):} \\
    اكتب الدوال الثابتة التالية في فئة البرنامج الرئيسي (بجوار الـ \textenglish{Main})، ثم استدعها لاختبارها:
    \begin{enumerate}[label=\textbf{\alph*)}]
        \item عمليّة باسم \textenglish{GetFurthestFromOrigin} تتلقّى مصفوفة من النقاط (مصفوفة كائنات \textenglish{Point}) وتُعيد النقطة التي بُعدها عن نقطة الأصل $(0,0)$ هو الأكبر.
        \item عمليّة باسم \textenglish{GetDistancesFromOrigin} تتلقّى مصفوفة من النقاط وتُعيد مصفوفة جديدة موازية، بحيث تحتوي كل خانة فيها على بُعد النقطة المقابلة عن نقطة الأصل $(0,0)$.
        \item عمليّة باسم \textenglish{GetMostRightPoint} تتلقى مصفوفة من النقاط. تعيد العملية مؤشرًا يُشير إلى النقطة الموجودة في أقصى اليمين (أي النقطة التي لها أكبر إحداثي $x$).
        \item عمليّة باسم \textenglish{FindClosestPair} تتلقى مصفوفة من النقاط. تقوم العملية بإيجاد أقرب نقطتين على بعضهما البعض في المصفوفة، وتطبع إحداثياتهما. افترض أن المصفوفة تحتوي على نقطتين على الأقل.
    \end{enumerate}

\end{enumerate}

\ifwithsols
{\small
\begin{boxSolution}[حلول ورقة العمل]
\begin{english}
\begin{minted}{csharp}
using System;

public class Point
{
    private double x;
    private double y;

    public Point(double x, double y)
    {
        this.x = x;
        this.y = y;
    }

    public Point()
    {
        this.x = 0;
        this.y = 0;
    }

    public Point(Point other)
    {
        this.x = other.x;
        this.y = other.y;
    }

    public double GetX() { return this.x; }
    public double GetY() { return this.y; }
    public void SetX(double x) { this.x = x; }
    public void SetY(double y) { this.y = y; }

    public void MoveBy(double dx, double dy)
    {
        this.x += dx;
        this.y += dy;
    }

    public string ToString()
    {
        return $"({this.x}, {this.y})";
    }

    // Static Methods in Point class
    public static double Slope(Point p1, Point p2)
    {
        return (p2.y - p1.y) / (p2.x - p1.x);
    }

    public static double Distance(Point p1, Point p2)
    {
        return Math.Sqrt(Math.Pow(p1.x - p2.x, 2) + Math.Pow(p1.y - p2.y, 2));
    }

    public static Point GetMiddlePoint(Point p1, Point p2)
    {
        return new Point((p1.x + p2.x) / 2, (p1.y + p2.y) / 2);
    }

}

public class Program
{
    // Static Methods in Main Program
    public static Point GetFurthestFromOrigin(Point[] points)
    {
        Point origin = new Point(0, 0);
        Point furthest = points[0];
        double maxDist = Point.Distance(furthest, origin);

        for (int i = 1; i < points.Length; i++)
        {
            double dist = Point.Distance(points[i], origin);
            if (dist > maxDist)
            {
                maxDist = dist;
                furthest = points[i];
            }
        }
        return furthest;
    }

    public static double[] GetDistancesFromOrigin(Point[] points)
    {
        Point origin = new Point(0, 0);
        double[] distances = new double[points.Length];
        for (int i = 0; i < points.Length; i++)
        {
            distances[i] = Point.Distance(points[i], origin);
        }
        return distances;
    }

    public static Point GetMostRightPoint(Point[] arr)
    {
        Point max = arr[0];
        for (int i = 1; i < arr.Length; i++)
        {
            if (arr[i].GetX() > max.GetX())
                max = arr[i];
        }
        return max;
    }

    public static void FindClosestPair(Point[] arr)
    {
        double minDistance = double.MaxValue;
        Point p1 = null;
        Point p2 = null;

        for (int i = 0; i < arr.Length; i++)
        {
            for (int j = i + 1; j < arr.Length; j++)
            {
                double dist = Point.Distance(arr[i], arr[j]);
                if (dist < minDistance)
                {
                    minDistance = dist;
                    p1 = arr[i];
                    p2 = arr[j];
                }
            }
        }
        if (p1 != null && p2 != null)
            Console.WriteLine($"{p1.ToString()} and {p2.ToString()}");
    }

    public static void Main(string[] args)
    {
        // Tests can be written here...
    }
}
\end{minted}
\end{english}
\end{boxSolution}
}
\fi

\end{document}
